\section{知识产权与自主可控}
\label{sec:ip}

\subsection{为什么 Apache 2.0 不等于"无壁垒"}

Vulca 选择 \keyword{Apache 2.0} 作为主仓库协议，有两个目的：
\keyword{降低 agent 开发者试用门槛}（无审批即可 pip 安装），
与\keyword{建立长期可信的开源品牌}（任何厂商可以审计源码，避免"封闭魔法"式不信任）。

但开源不等于没有壁垒。Vulca 的真实护城河建立在\keyword{三类可持续积累的资产}上，
这三类都不是靠一份代码协议能复刻的：

\begin{center}
\begin{tabularx}{\linewidth}{@{}>{\tabRR}p{3.4cm}>{\tabRR}X@{}}
\toprule
\rowcolor{tableHead}
\heiti 资产类别 & \heiti 具体构成 \\
\midrule
\keyword{数据资产} & 13 种艺术/设计传统的 L1–L5 YAML 评价体系；社区扩展贡献的新传统；行业客户定制包 \\
\keyword{工程资产} & MCP 工具标准 I/O schema、skill 模板、高覆盖率测试套件、多 provider 抽象层、本地化链路工程 \\
\keyword{生态资产} & Claude Code 插件市场位置、GitHub 贡献者网络、学术合作（ACM MM / EmoArt）、模型厂商直联 credits \\
\bottomrule
\end{tabularx}
\end{center}

\subsection{开源协议策略（边界清晰）}

\begin{itemize}
  \item \keyword{核心 SDK（Apache 2.0）}：\code{vulca} 主包、MCP 工具集、13 种基础文化传统 YAML、provider 抽象层、skill 模板；
  \item \keyword{商业专有部分}：Vulca Cloud 后端（托管、SSO、账单、审计日志）、行业客户定制 YAML 包、
  On-Prem 打包脚本、企业级监控与升级支持；
  \item \keyword{CLA（贡献者协议）}：对第三方提交到主仓库的代码或 YAML 要求签署 CLA，\keyword{允许将来的商业化组合使用}；
  \item \keyword{品牌商标}：\code{Vulca} 名称与 logo 建议\keyword{国内第 9 类（软件）、第 41 类（教育培训）、第 42 类（SaaS/设计服务）}全部注册，条件允许时追加 Madrid 国际注册 \presumed{建议默认全选三个类}。
\end{itemize}

\subsection{自主可控证明链}

针对 C 类受众（孵化器 / 政府立项 / 高校转化）最关注的\keyword{自主可控}问题，
Vulca 有一条可以被审查的证明链：

\begin{enumerate}
  \item \keyword{代码自主可控}\quad 主包 Apache 2.0 开源，\keyword{零闭源依赖}运行（除了可选的云端 provider）；
  \item \keyword{算法自主可控}\quad 核心调度逻辑、文化评估规则、分层编排算法都写在 Vulca 自己的代码里，\keyword{不依赖任何厂商 SDK}；
  \item \keyword{模型可替换}\quad provider 抽象层支持 ComfyUI（本地开源）/ Gemini / OpenAI / mock 四条路径；
  \item \keyword{数据自主可控}\quad 本地链路所有图像、评估结果、会话历史都留存在客户本地磁盘（\code{artifacts/})，\keyword{不向任何云服务上报}；
  \item \keyword{算力自主可控}\quad Apple Silicon MPS 已端到端验证；CUDA 适配在 2026 Q4 前完成；国产算力首轮适配路线：\keyword{2027 Q1 完成 Ascend POC，2027 Q3 完成寒武纪适配} \presumed{与国产算力政策窗口对齐的默认时间表}；
  \item \keyword{合规自主可控}\quad 下表给出主要政策与 Vulca 设计响应的\keyword{条款级}对应（示意口径，正式申报时依评审口径微调）：
\end{enumerate}

\paragraph{政策条款 → Vulca 设计响应（示意映射）}
\begin{center}
\small
\begin{tabularx}{\linewidth}{@{}>{\tabRR}p{5.0cm}>{\tabRR}X@{}}
\toprule
\rowcolor{tableHead}
\heiti 政策 / 条款 & \heiti Vulca 设计响应 \\
\midrule
《生成式 AI 服务管理办法》第 8 条（训练数据来源合规） & Vulca 自身不训练底层模型；上游 provider 可替换；本地路径允许客户完全使用自有模型 \\
《生成式 AI 服务管理办法》第 11 条（算法备案） & 默认判定：\keyword{Cloud 版本（面向公众服务）需备案}；\keyword{SDK + On-Prem（不面向公众）不需备案} \presumed{正式申报前由法务顾问最终判定} \\
《生成式 AI 服务管理办法》第 14 条（内容审核义务） & 关键词黑名单 + 上游 provider 审核回调；敏感类别检测；违规内容日志；Cloud 模式下拒绝生成并记录 \\
《数据出境安全评估办法》 & 本地 provider 路径\keyword{原生满足不出境}；Cloud 路径可选择国内节点 provider 或强制本地 provider \\
《"十四五"数字经济发展规划》文化数字化章节 & Vulca 文化评估 L1–L5 YAML 对应"传统文化数字化复现"核心能力；可开放给国家文博数字化项目免费使用 \\
\textit{著作权 / 开源协议} & 主仓库 Apache 2.0 对外；商业组件单独许可；第三方贡献走 CLA，避免协议混合 \\
\bottomrule
\end{tabularx}
\end{center}

\subsection{知识产权规划}

\subsubsection*{已形成的 IP 资产}
\begin{itemize}
  \item \keyword{开源著作权}\quad 主仓库 commit history 完整，著作权归 Vulca 团队；
  \item \keyword{文化评估体系}\quad 13 种传统 YAML 本身可作为\keyword{数据库著作权}声明（YAML 结构与内容的原创性聚合）；
  \item \keyword{技术文档体系}\quad \code{docs/superpowers/specs/} 下的设计文档、\code{docs/} 下的工程手册，构成\keyword{独立的技术资产}。
\end{itemize}

\subsubsection*{计划申请的 IP}

\presumed{默认 18 个月内 IP 申请计划：以下 4 类全选；总成本约 ¥15–25 万，与 §8.2 法务预算对齐}

\begin{itemize}
  \item \keyword{软件著作权登记}\quad 主仓库 + 文化评估体系各 1 件，成本低、周期短，建议 T+3 月内完成；
  \item \keyword{发明专利}\quad 候选：\keyword{agent-native MCP 工具编排方法}、\keyword{基于 L1–L5 分层评估的多文化图像审美判别方法}、\keyword{多 provider 图像工作流的动态降级策略}。建议 2 件，T+6 月前完成提交；
  \item \keyword{商标}\quad \code{Vulca} 文字 + 设计标，第 9 类（软件）、第 42 类（SaaS/设计服务）、第 41 类（教育培训）；
  \item \keyword{美术作品著作权}\quad logo、官网视觉、核心 diagram 登记。
\end{itemize}

\subsection{社会效益与可量化指标}
\label{sec:ip:social}

作为 C 类受众（政府立项 / 孵化器 / 高校转化）的\keyword{打分维度}之一，
Vulca 在社会效益上提供\keyword{可年度核对的量化指标}，而不是口号式描述。
以下目标值为\keyword{12–24 个月内}可达成的保守口径：

\begin{center}
\small
\begin{tabularx}{\linewidth}{@{}>{\tabRR}p{4.2cm}>{\tabRR}X>{\tabRR}p{2.6cm}@{}}
\toprule
\rowcolor{tableHead}
\heiti 指标 & \heiti 说明 & \heiti 12-24 月目标 \\
\midrule
\keyword{服务机构数} & 博物馆 / 出版社 / 品牌 / 教育机构使用 Vulca（免费社区或商业部署均计入） & $\ge$ 20 家 \presumed{} \\
\keyword{数字化作品数} & 通过 Vulca 分层 / 编辑 / 评估的文化作品累计数 & $\ge$ 5000 件 \presumed{} \\
\keyword{非遗 / 馆藏 POC 案例} & 公开发表的非遗复现、馆藏数字化 POC 项目数 & $\ge$ 3 个 \presumed{} \\
\keyword{学术 / 教育课时} & 高校艺术 / 设计教学中使用 Vulca 的课时数（合作院校披露） & $\ge$ 200 课时 \presumed{} \\
\keyword{社区贡献 YAML} & 开源社区贡献的新文化传统定义数 & $\ge$ 3 个（跨机构）\presumed{} \\
\keyword{就业拉动} & 团队直接雇佣 + 生态合作方兼职 & $\ge$ 8 人全职 + 10 人兼职 \presumed{} \\
\bottomrule
\end{tabularx}
\end{center}

每个指标配合\keyword{公开可核验的渠道}：GitHub / PyPI / 合作机构新闻稿 / 学术论文致谢 / 年度报告。
申报阶段可据此加权生成对申报单位的\keyword{社会效益承诺书}，纳入里程碑 KPI。

\subsection{数据合规与伦理}

\begin{itemize}
  \item \keyword{训练数据使用}\quad Vulca 自身\keyword{不训练底层生成模型}，所有生成能力通过 provider 调用既有模型（SDXL、FLUX、Gemini 等），训练数据责任由上游模型提供方承担；
  \item \keyword{文化评估数据}\quad YAML 定义基于公开艺术史资料撰写；建议为每一条规则注明引用来源，以支持学术与政策层面的可追溯；
  \item \keyword{用户数据}\quad 本地模式下不上传；Cloud 模式下遵循最小化收集原则，用户可一键导出和删除；
  \item \keyword{伦理审查}\quad 文化评估不做"高低优劣"判断，仅做"是否符合 X 传统 L1–L5 描述"的判别，\keyword{避免文化偏见}写入产品；
  \item \keyword{未成年人 / 敏感内容}\quad 默认安全策略：\keyword{三层防线}——(i) 输入端关键词黑名单拦截；(ii) 上游 provider 自带审核回调；(iii) 输出端二分类敏感检测（默认模型为 CLIP-based 开源检测器）；未成年人场景提供"教学白名单"模式。所有拦截事件进入审计日志。\presumed{正式产品上线前完成安全评估与备案}
\end{itemize}
